% ============================================================
%  7.1. COMPRENSIÓN DEL NEGOCIO  -- parte del capítulo 7
%  ------------------------------------------------------------
%  Guía 2.7.1: (1) contexto operativo o institucional,
%  (2) traducción de la necesidad, (3) resultado esperado y
%  criterio de utilidad, (4) enlace con lo ya escrito.
%  No es "Comprensión de los datos" (fuentes, variables, EDA):
%  eso va en 08_02.
%  Fuente de verdad: trufi-data-science (EXPLICATIVO.md,
%  config.py, resultados/<fase>/metricas.json).
% ============================================================
\section{Comprensión del negocio}
\label{sec:comprension-negocio}

Esta primera fase de CRISP-DM define para quién se trabaja, qué decisión se quiere apoyar y
cómo se reconocerá que el resultado sirve. No trata todavía de los datos, que se describen
en la sección~\ref{sec:comprension-datos}, sino del proceso real de la organización y del
modo en que su necesidad se convierte en un problema que la Ciencia de Datos puede
abordar.

\subsection{Contexto operativo e institucional}
\label{sec:contexto-operativo}

Trufi Association es una organización sin fines de lucro que mantiene Trufi App, un
planificador de viajes para el transporte público del eje metropolitano de Cochabamba.
Como la mayor parte de ese transporte opera bajo la modalidad de paratránsito, sin
horarios ni paradas oficiales publicadas (sección~\ref{subsec:paratransito}), la
información que ofrece la aplicación no proviene de los operadores, sino de un
relevamiento colaborativo. El proceso real tiene cuatro pasos:
\begin{enumerate}
  \item Un grupo de voluntarios recorre una línea de micro o trufi.
  \item Los voluntarios registran su trayecto en OpenStreetMap
        (sección~\ref{subsec:cartografia-colaborativa}).
  \item El trazado se verifica con los operadores o en terreno.
  \item El trazado se convierte al formato GTFS (sección~\ref{subsec:gtfs}), que es el que
        consume el planificador de la aplicación.
\end{enumerate}
El feed resultante, empleado en este trabajo, describe 141 líneas de 43 operadores.

Cada jornada de mapeo consume desplazamientos, tiempo de edición y coordinación entre
personas que colaboran de forma voluntaria. Por eso la capacidad de la organización es
limitada, y cada jornada debe dirigirse a alguna zona concreta del territorio. Hoy esa
elección descansa en el conocimiento acumulado del equipo y en observaciones informales.
Al mismo tiempo, la aplicación genera de forma pasiva un registro de cada consulta de ruta
que realizan sus usuarios, con origen, destino, hora e identificador anónimo de
instalación. Ese registro abarca casi dos años (septiembre de 2022 a junio de 2024) y más
de 1,9~millones de consultas, pero no forma parte del proceso con el que se decide dónde
mapear. El punto de partida del proyecto es esa desconexión entre un proceso operativo que
necesita priorizar y un registro de uso que no se aprovecha para hacerlo.

\subsection{Traducción de la necesidad a un problema de Ciencia de Datos}
\label{sec:traduccion-necesidad}

La necesidad de la organización se expresa en términos operativos: \emph{¿a qué zonas
conviene enviar primero a los voluntarios?} Esa pregunta no puede responderse contando
consultas, porque un conteo bajo es ambiguo (sección~\ref{sec:descripcion-problema}).
Puede tratarse de una zona con poca población, o de una zona donde la aplicación no ofrece
información útil y, por eso, casi nadie la consulta. Para separar ambas situaciones hace
falta una referencia: cuántas consultas \emph{cabría esperar} en cada zona si se
comportara como las zonas comparables.

Con esa idea, la necesidad se traduce en un problema analítico de dos pasos:

\begin{enumerate}
  \item \textbf{Estimar el conteo esperado de consultas por zona.} Es un problema de
    regresión sobre datos de conteo (sección~\ref{subsec:datos-conteo}). La variable a
    estimar es el total de consultas originadas en cada zona H3 de resolución~8 durante el
    período, y la población residente actúa como \emph{exposición}: a igualdad de
    condiciones, una zona con el doble de habitantes debería generar el doble de consultas.
    Lo esperado se estima con información que no depende de la red mapeada: la población y
    la localización de la zona.
  \item \textbf{Comparar lo observado con lo esperado.} La diferencia entre ambos,
    expresada en una escala comparable entre zonas grandes y pequeñas, se denomina
    \emph{brecha} (sección~\ref{subsec:observado-esperado}). Una brecha negativa indica que
    la zona consulta menos de lo que cabría esperar por su población y su entorno. Esa
    señal, contrastada después con la cobertura de la red mapeada, orienta la revisión en
    terreno.
\end{enumerate}

Conviene precisar lo que el problema \emph{no} es. No es un pronóstico temporal, porque no
se predice cuántas consultas habrá la semana próxima, sino cuántas corresponden a cada zona
en el período observado. Tampoco es una clasificación de zonas en ``bien'' o ``mal''
mapeadas, porque no existe una etiqueta verificada de cobertura real con la cual entrenar
un clasificador. La cobertura de la red mapeada se excluye deliberadamente del cálculo de
lo esperado: si se incluyera, el modelo aprendería que las zonas sin rutas consultan poco y
les asignaría una expectativa baja, con lo que la brecha que se busca detectar desaparecería
por construcción.

\subsection{Resultado esperado y criterio de utilidad}
\label{sec:criterio-utilidad}

El resultado que se espera entregar a Trufi Association no es ``un modelo'', sino tres
productos que puedan usarse sin conocimientos de programación:
\begin{itemize}
  \item un archivo con el valor observado, el esperado y la brecha de cada zona;
  \item un mapa interactivo de la brecha;
  \item una lista corta de veinte zonas priorizadas para la revisión en terreno.
\end{itemize}
La mejora que se persigue es operativa: sustituir la elección informal entre más de mil
zonas por una lista corta, ordenada y justificada con datos.

Para que esa lista sea confiable, la técnica que produce el esperado debe cumplir criterios
medibles, fijados antes de ajustar cualquier modelo. La técnica se elegirá entre siete
alternativas de complejidad creciente. Sus códigos (B0 a B3 y M1 a M3) se explican en la
sección~\ref{sec:seleccion-algoritmos}. La tabla~\ref{tab:criterios-utilidad} resume los
criterios, cómo se mide cada uno y qué pregunta de la organización responde.

\begin{table}[H]
  \centering
  \caption{Criterios de utilidad declarados antes de modelar}
  \label{tab:criterios-utilidad}
  \interlineadoSimple
  \small
  \begin{tabularx}{\textwidth}{@{}L{0.23\textwidth}L{0.33\textwidth}Y@{}}
    \toprule
    \textbf{Criterio} & \textbf{Cómo se mide} & \textbf{Pregunta que responde} \\
    \midrule
    Supera al punto de referencia inicial & $D^2$ frente a la tasa global en zonas no vistas, mayor que cero & ¿El contexto territorial aporta algo más que la población? \\
    Ordena bien las zonas & Correlación de Spearman entre observado y esperado & ¿Las zonas que más consultan quedan arriba? \\
    Está calibrada & Suma de lo esperado dividida entre suma de lo observado, cercana a 1 & ¿El volumen total estimado es realista? \\
    Es estable & Índice de Jaccard del top-20 ante variaciones razonables & ¿La lista cambia si se modifica un supuesto menor? \\
    Es accionable & Número de zonas que el equipo debe revisar & ¿Se reduce el trabajo de decisión? \\
    \bottomrule
  \end{tabularx}
  \fuente{Elaboración propia. Las métricas se definen en la sección~\ref{sec:metricas}; los
  valores obtenidos se presentan en la sección~\ref{sec:evaluacion-resultados}.}
\end{table}

A estos criterios se suman dos reglas de protocolo, también declaradas de antemano. La
primera es que la técnica se elige con una regla escrita antes de ver los resultados, y no
por ser la primera que se probó ni la de menor error aparente
(sección~\ref{sec:regla-seleccion}). La segunda es que el conjunto de prueba se evalúa una
sola vez, al final.

La señal definitiva de utilidad, sin embargo, no es estadística. Sería la proporción de
zonas visitadas en las que los voluntarios encuentran rutas sin mapear. Esa verificación
en terreno excede el alcance de la materia (capítulo~\ref{cap:alcance}) y no se ha
realizado. Para que pueda medirse en cuanto empiecen las visitas, el archivo
\var{zonas_prioritarias.csv} incluye dos columnas de registro:
\begin{itemize}
  \item \var{field_result}, que registra el resultado de la visita con una de cuatro
        categorías: \var{ruta_no_mapeada}, \var{ruta_mapeada_correcta}, \var{sin_transporte}
        o \var{sin_visitar}. Esta última viene precargada en las veinte zonas.
  \item \var{field_notes}, que queda libre para observaciones.
\end{itemize}
El protocolo correspondiente se describe en la sección~\ref{sec:productos-despliegue}.

\subsection{Enlace con el problema, el alcance y los objetivos}
\label{sec:enlace-objetivos}

Esta fase concreta en términos operativos lo que los capítulos anteriores plantearon en
términos de investigación:
\begin{itemize}
  \item La ambigüedad del conteo descrita en el capítulo~\ref{cap:problema} se resuelve con
        la comparación entre lo observado y lo esperado.
  \item Las exclusiones del capítulo~\ref{cap:alcance} (sin pronóstico temporal, sin datos
        externos ni imputación de semanas faltantes, sin despliegue en producción) delimitan
        las técnicas admisibles y el tipo de entrega.
  \item Cada objetivo específico del capítulo~\ref{cap:objetivos} corresponde a una o dos
        fases de CRISP-DM y deja una evidencia verificable en el repositorio del proyecto,
        como se muestra en la tabla~\ref{tab:trazabilidad-oe}.
\end{itemize}

\begin{table}[H]
  \centering
  \caption{Correspondencia entre objetivos específicos, fases y evidencias}
  \label{tab:trazabilidad-oe}
  \interlineadoSimple
  \small
  \begin{tabularx}{\textwidth}{@{}L{0.07\textwidth}L{0.25\textwidth}L{0.17\textwidth}Y@{}}
    \toprule
    \textbf{OE} & \textbf{Fase y notebook} & \textbf{Sección} & \textbf{Evidencia principal} \\
    \midrule
    OE1 & Comprensión de los datos (\var{01_eda}); preparación (\var{02_preprocesamiento}, \var{03_feature_engineering}) & \ref{sec:comprension-datos} y \ref{sec:preparacion-datos} & Tabla de 1.628 zonas con consultas, población y cobertura; \var{metricas.json} de las tres fases \\
    OE2 & Preparación (reserva de prueba, \var{03_feature_engineering}); modelado (\var{04_modelado}) & \ref{sec:conjuntos-entrenamiento-prueba} y \ref{sec:modelado} & Bloques de prueba (\var{test_blocks.csv}); comparación de siete técnicas (\var{cv_resumen.csv}); regla de selección (\var{regla_seleccion.csv}) \\
    OE3 & Evaluación (\var{05_evaluacion}) & \ref{sec:evaluacion-resultados} & Métricas de prueba, brecha por zona, contraste con la cobertura y sensibilidad \\
    OE4 & Despliegue como propuesta (\var{06_propuesta}) & \ref{sec:despliegue} & Predicciones por zona, lista de veinte zonas prioritarias, mapa interactivo y procedimiento de actualización \\
    \bottomrule
  \end{tabularx}
  \fuente{Elaboración propia, a partir de la estructura del repositorio
  \var{trufi-data-science} (carpetas \var{notebooks/} y \var{resultados/}).}
\end{table}

La figura~\ref{fig:flujo-crisp-proyecto} muestra cómo se instanciaron las seis fases de
CRISP-DM en el proyecto. Cada fase de datos corresponde a un notebook que lee lo que dejó
la anterior, regenera sus salidas y termina con verificaciones automáticas de sus propias
cifras.

El carácter iterativo de la metodología (sección~\ref{sec:crisp-dm}) se manifestó en dos
retornos concretos a esta fase, señalados con flechas discontinuas en la figura:
\begin{itemize}
  \item El análisis exploratorio fijó la delimitación operativa del área de estudio, que
        el capítulo~\ref{cap:alcance} había dejado abierta a propósito
        (sección~\ref{sec:area-estudio}).
  \item La evaluación mostró que la lista de zonas con mayor brecha estaba formada casi por
        completo por zonas ya cubiertas por la red mapeada
        (sección~\ref{sec:interpretacion-resultados}). Como en esas zonas mapear no aporta,
        se precisó el criterio de utilidad: la lista priorizada se restringe a zonas sin
        trazado cercano.
\end{itemize}
La técnica elegida no se modificó por ese cambio. Ambas decisiones quedan documentadas
como tales, y no se presentan como parte del diseño original.

\begin{figure}[H]
  \centering
  \caption{Flujo CRISP-DM del proyecto: fases, notebooks y objetivos específicos}
  \label{fig:flujo-crisp-proyecto}
  \begin{tikzpicture}[
      font=\small,
      fase/.style={draw, rounded corners=3pt, fill=gray!8, align=center,
                   text width=4.0cm, minimum height=2.1cm, inner sep=3pt},
      flecha/.style={-{Stealth[length=2.5mm]}, thick},
      retorno/.style={-{Stealth[length=2.5mm]}, thick, dashed, gray!70!black},
      node distance=0.9cm and 0.45cm]
    \node[fase] (f1) {\textbf{1. Comprensión\\del negocio}\\[2pt] sección~\ref{sec:comprension-negocio}};
    \node[fase, right=of f1] (f2) {\textbf{2. Comprensión\\de los datos}\\[2pt] {\footnotesize\texttt{01\_eda}}\\ OE1};
    \node[fase, right=of f2, text width=5.0cm] (f3) {\textbf{3. Preparación\\de los datos}\\[2pt] {\footnotesize\texttt{02\_preprocesamiento}}\\ {\footnotesize\texttt{03\_feature\_engineering}}\\ OE1 y OE2};
    \node[fase, below=1.4cm of f3, text width=5.0cm] (f4) {\textbf{4. Modelado}\\[2pt] {\footnotesize\texttt{04\_modelado}}\\ OE2};
    \node[fase, left=of f4] (f5) {\textbf{5. Evaluación}\\[2pt] {\footnotesize\texttt{05\_evaluacion}}\\ OE3};
    \node[fase, left=of f5] (f6) {\textbf{6. Despliegue}\\ (propuesta)\\[2pt] {\footnotesize\texttt{06\_propuesta}}\\ OE4};
    \draw[flecha] (f1) -- (f2);
    \draw[flecha] (f2) -- (f3);
    \draw[flecha] (f3) -- (f4);
    \draw[flecha] (f4) -- (f5);
    \draw[flecha] (f5) -- (f6);
    \draw[retorno] (f2.north) to[bend right=25]
      node[above, font=\footnotesize, text=black] {delimitación del área} (f1.north);
    \draw[retorno] ($(f5.north west)+(0.5,0)$) --
      node[right, pos=0.5, font=\footnotesize, text=black] {regla de priorización}
      ($(f1.south east)+(-0.5,0)$);
  \end{tikzpicture}
  \fuente{Elaboración propia. Las flechas discontinuas señalan los dos retornos documentados
  a la comprensión del negocio.}
\end{figure}
